\originalchapter{连续自旋、涡旋与BKT转变}
\label{ch:bkt}
\source{8.334 L20--L22。}
\originalsection{非线性sigma模型}
$O(n)$低温固定幅度场$n(x)$满足$n^2=1$，作用$\mathcal H=(1/2g)\int(\nabla n)^2$，$g$与温度/刚度之比有关。参数化$n=(\pi,\sqrt{1-\pi^2})$，展开得$\mathcal H=(1/2g)\int[(\nabla\pi)^2+(\pi\cdot\nabla\pi)^2+\cdots]$，约束造成梯度相互作用，不能把$n$的$n$个分量都当独立Gaussian。

在$d=2+\epsilon$，低温动量壳积分的自旋波重整化给$\dd g/\dd l=-\epsilon g+(n-2)g^2/(2\pi)+\cdots$。$n>2,d=2$耦合随尺度增长，相关长约$a\exp[2\pi/((n-2)g)]$，自旋波弱耦合最终失控。这个一回路系数来自约束相互作用与场重标的组合，$n=2$时其抵消与可用单角度Gaussian表述一致；但是涡旋仍是非微扰自由度。
\originalsection{XY自旋波与代数序}
$n=2$写$n=(\cos\theta,\sin\theta)$，$\mathcal H=K\int(\nabla\theta)^2/2$，$K=\rho_s/(\kb T)$。在二维高斯场$\avg{[\theta(r)-\theta(0)]^2}=\log(r/a)/(\pi K)$，所以$\avg{\e^{\ii\theta(r)}\e^{-\ii\theta(0)}}=(a/r)^{\eta}$，$\eta=1/(2\pi K)$。相位差的对数增长使通常磁化为0，却有幂律关联，正是准长程序。

紧致角还允许$\oint\nabla\theta\cdot\dd l=2\pi m$涡旋。单涡$\theta=m\arg(x+\ii y)$，梯度$m/r$，能量$E_v=\pi\rho_sm^2\log(L/a)+E_c$。可放置位置的熵约$2\kb\log(L/a)$，对单位涡的简单能熵竞争给临界刚度量级$\rho_s/(\kb T)=2/\pi$。因其他涡与自旋波会重整刚度，应使用长距离的$K_R$，不能把这个估计中的裸刚度直接叫精确转变温度。
\originalsection{Coulomb气与介电屏蔽}
多个涡的能量在中性条件$\sum_im_i=0$下为$-2\pi\rho_s\sum_{i<j}m_im_j\log(r_{ij}/a)+E_c\sum_im_i^2$。小涡反涡偶极子在大涡产生的场中极化，屏蔽对数作用、降低有效$K$。赋予单位涡逸度$y=\e^{-E_c/\kb T}$，积分尺寸$a$到$a\e^{\dd l}$的偶极子并重标，标准归一给
\begin{equation}
\frac{\dd K^{-1}}{\dd l}=4\pi^3y^2+O(y^4),\qquad\frac{\dd y}{\dd l}=(2-\pi K)y+O(y^3).
\end{equation}
第二式的2是位置积分维数，$-\pi K$是涡自能随截止变动的尺度代价；第一式正号表明偶极屏蔽只降低刚度，系数依$y$的定义。保持中性与偶极角平均才能获得这个二阶修正。
\originalsection{流、普适跳跃与本质奇异}
$y=0,K>2/\pi$是一条低温固定线，涡无关，相关幂随$K_R$连续变化；$K<2/\pi$时$y$增长，自由涡破坏代数序。分界流抵达$K_R=2/\pi$，故$\rho_s(T_c^-)=2\kb T_c/\pi$，临界$\eta=1/4$。

在端点附近令$x=2-\pi K$并适当重标$Y\propto y$，最低阶$\dot x\propto Y^2,\dot Y=xY$，有双曲不变量$x^2-Y^2$。高温距分界$\tau>0$时流在临界区域停留$l_*\propto1/\sqrt\tau$，故$\xi\sim a\exp(b/\sqrt\tau)$。这是本质奇异，不能赋一个普通有限$\nu$来模拟。$b$依微观逸度，刚度跳跃与临界指数则由固定点结构决定。

有限系统即使无穷体系无磁化也可有大$\avg{|M|}$；应测关联与扭曲刚度随$L$的对数修正，再区分BKT与普通幂律临界。孤立单涡不能在严格周期环面中净存在，模拟要用中性涡对或相应拓扑扇区。
